% !PRINTED 275
\weberpage{275}
Hierin bedeutet $\gamma$ eine primitive Wurzel, also $\gamma^{\frac{p^2-1}{\nu}}$
eine primitive $\nu^{\mathrm{te}}$ Einheitswurzel des Congruenzkörpers
$\mathfrak{E}_{2,p}$ (\S.\,67).

3. \gesperrt{Beide Pole einer Substitution $\Theta$ sind gleichzählig}
(\S.\,52, 4.).

Ist $Q$ die cyklische Gruppe $\nu^{\mathrm{ten}}$ Grades der Potenzen
von $\Theta$ und $n = \nu\mu$, so erhält man
\[
G = Q + \psi_1 Q + \cdots + \psi_{\mu-1} Q,
\]
worin $\psi_1, \ldots, \psi_{\mu-1}$ Substitutionen aus $G$ sind, und

4. \gesperrt{die Zahlen}
\weq{4}{\alpha,\ \psi_1(\alpha) = \alpha_1,\ \psi_2(\alpha) = \alpha_2,\ \ldots,\ \psi_{\mu-1}(\alpha) = \alpha_{\mu-1}}
\gesperrt{bilden ein System gleichzähliger conjugirter Pole der Gruppe $G$.}

Die sämmtlichen Pole der Gruppe $G$ lassen sich also in Systeme
conjugirter Pole zusammenfassen, und wir bekommen so die unbestimmte
Gleichung
\weq{5}{\begin{aligned}
2n - 2 &= \mu(\nu-1) + \mu'(\nu'-1) + \mu''(\nu''-1) \ldots \\
&= nh - \mu - \mu' - \mu'' \ldots,
\end{aligned}}
wenn $h$ die Anzahl der Systeme conjugirter Pole ist, und von dieser
Gleichung haben wir in \S.\,52 gesehen, dass sie nur eine beschränkte
Anzahl von Lösungen zulässt.

Um die diesen Lösungen entsprechenden Theiler der Congruenzgruppen zu
finden, können wir geradezu in den Formeln der Polyëdergruppen, die im
achten Abschnitte aufgestellt sind, für die darin vorkommenden
Einheitswurzeln die in \S.\,67 definirten Einheitswurzeln des Körpers
$\mathfrak{E}_{2,p}$ setzen, mit denen ja nach denselben Regeln
gerechnet wird. Dabei können natürlich nur solche Einheitswurzeln
vorkommen, deren Grad ein Theiler von $p^2 - 1$ ist.

Es ist schliesslich bei jeder solchen Gruppe noch zu untersuchen, ob
ihre Substitutionen in $L_p$ oder in einer damit isomorphen Gruppe
enthalten sind. Wir haben dann folgende Fälle von Lösungen der
Gleichung (5), wobei $\gamma$ stets eine primitive Wurzel des Körpers
$\mathfrak{E}_{2,p}$ bedeutet (\S.\,55).

\gesperrt{I. Cyklische Gruppen vom Grade $n$,}\quad
$\begin{aligned} h &= 2,\ \nu = \nu' = n \\ \mu &= \mu' = 1 \end{aligned}$
\[
C_n = \begin{pmatrix} \gamma^{\frac{p^2-1}{2n}\lambda}, & 0 \\[2pt] 0, & \gamma^{-\frac{(p^2-1)}{2n}\lambda} \end{pmatrix},
\quad \lambda = 0, 1, \ldots, n-1.
\]
